\section{Implementação}
\label{sec:implementacao}

\subsection{Onde o código vive}

\begin{center}
\begin{tabular}{ll}
  \toprule
  \texttt{higflow/src/hig-flow-fronteira-imersa.\{h,c\}} & o núcleo do método \\
  \texttt{higflow/examples-common/fronteira-imersa.c}    & adaptador solver/módulo \\
  \texttt{higflow/fronteira-imersa/testes/}              & verificação isolada \\
  \texttt{higflow/example2d\_SchaeferTurek/}             & o experimento \\
  \bottomrule
\end{tabular}
\end{center}

\paragraph{O módulo não conhece o solver, e isso é deliberado.} Ele recebe um
\texttt{sim\_facet\_domain} e uma \texttt{distributed\_property} por direção, não
um \texttt{higflow\_solver}. Foi essa separação que permitiu verificar a
conservação sem montar um solver (Seção~\ref{sec:verificacao}).

\paragraph{O acoplamento é por ponteiro de função.} O arquivo
\texttt{hig-flow-step.c} é ligado por \emph{todos} os exemplos. Uma chamada
direta obrigaria os doze \emph{Makefiles} a ligar o módulo por um recurso que um
exemplo usa. Com o gancho, quem não instala corpo não referencia símbolo algum.
A primeira tentativa, com chamada direta, levou a suíte de 33 para 0 de 33 casos.

\subsection{Estrutura de dados da malha lagrangeana}

Os marcadores vivem num \texttt{DMSwarm} do PETSc --- biblioteca já presente no
projeto, sem dependência nova. A entrada é uma \textbf{curva fechada por
segmentos}; em 3D, a extrusão dessa curva, formando uma superfície.

\paragraph{Marcadores nos centros dos subsegmentos, não nos vértices.} Marcador
em vértice duplicaria o vértice compartilhado entre segmentos consecutivos, e a
soma dos pesos deixaria de ser o perímetro. Com centros, o peso de cada marcador
é o comprimento do seu subsegmento, e a soma dá o perímetro exato --- invariante
que o teste cobra com tolerância de $10^{-12}$.

\paragraph{Distribuição por posse euleriana.} Cada marcador pertence ao rank que
possui a célula que o contém, condição para o espalhamento ser local. Uma guarda
aborta se a soma dos marcadores reivindicados não bater com o total: marcador
sobre fronteira de partição reivindicado por dois ranks (ou por nenhum) daria
força errada sem sintoma visível.

\subsection{O suporte do marcador, sem varredura}

O núcleo tem suporte de $1{,}5$ células para cada lado. Localizar as facetas do
suporte varrendo seria $O(N_L \times N_{\mathrm{facetas}})$. Não é preciso: a
malha é uniforme e os centros do suporte estão em \emph{deslocamentos conhecidos}
a partir de uma âncora, cada um achado por uma localização de ponto.

\paragraph{A âncora sai da célula, não de uma faceta.} A função
\texttt{sfd\_get\_facet\_with\_point} acha a faceta cujo \emph{plano} contém o
ponto; um marcador no meio de uma célula não está sobre plano algum naquela
direção. Ancorar numa faceta funcionava apenas para a direção em que o marcador
por acaso caísse sobre um plano de face --- e numa curva alinhada com a malha
isso atinge metade dos marcadores, com a força saindo exatamente pela metade.

\subsection{Acumulação franja \texorpdfstring{$\to$}{->} dono}

O espalhamento de um marcador próximo da borda escreve em facetas de \emph{outro}
rank. A sincronização da HiGTree (\texttt{dp\_sync}) envia dono\,$\to$\,franja e
\textbf{sobrescreve}: contribuições escritas numa franja seriam descartadas em
silêncio. A direção que falta é feita com \texttt{PetscSFReduce} e
\texttt{MPI\_SUM}, sobre um grafo montado a partir do que a HiGTree já tem
(\texttt{gid\_map}, \texttt{local\_count}, \texttt{total\_count}).

\paragraph{Contrato: o espalhamento é dono do campo que recebe.} Ele acumula
franja\,$\to$\,dono, o que só é correto se o campo contiver \emph{apenas}
contribuições espalhadas. Passar um campo já povoado faz o \emph{reduce} somar
um campo inteiro nas fronteiras de partição. Para somar num campo povoado ---
como $\mathbf{u}^*$ na rota pós-preditor --- espalha-se num campo zerado próprio
e soma-se depois.
